\section{Reproducibility Checklist}
\label{app:checklist}

We answer the Machine Learning Reproducibility Checklist v2.0 \citep{pineau} for the protocol at this commit. Items marked \ph{TODO} are pending the runs. Until the result files exist, every "yes" below describes the harness and the preregistered design, and the results section reports the protocol rather than a measurement.

\paragraph{Models and algorithms.}
\begin{itemize}
  \item Clear description of the mathematical setting: yes, \S\ref{sec:problem}.
  \item Clear explanation of assumptions: yes, \S\ref{sec:problem} and \S\ref{sec:limitations}.
  \item Complexity analysis: partial; wall-clock cost by tactic is reported in \S\ref{sec:setup}.
\end{itemize}

\paragraph{Datasets.}
\begin{itemize}
  \item Statistics and splits: yes, Table~\ref{tab:data} and \S\ref{sec:setup}.
  \item Exclusions and preprocessing: yes, \S\ref{sec:setup}.
  \item Downloadable data: \ph{link}, pending the prompt manifest.
  \item Collection process for new data: self-authored prompts from a documented template, plus the MarkMyWords tasks loaded from their release (Appendix~\ref{app:prompts}).
\end{itemize}

\paragraph{Shared code.}
\begin{itemize}
  \item Dependency specification: yes, \texttt{pyproject.toml} with version ranges.
  \item Evaluation code: partial; the statistics, split, strategy parser, and frontier helpers are in \texttt{efface/evaluate/} and \texttt{efface/search/}. Tactic modules, the detector adapter, and the three-pass search land with the first results.
  \item Pretrained models: named in Table~\ref{tab:attackers}; the open-weight generator checkpoints are named, and revisions are recorded in each run manifest.
  \item README with a table of results and the command that produces each: partial; the README carries the reproduction table with commands marked as planned.
\end{itemize}

\paragraph{Experimental results.}
\begin{itemize}
  \item Hyperparameter ranges and selection method: yes, Appendix~\ref{app:hyper}; the operating intensity is chosen by the pass-1 rule in \S\ref{sec:search}.
  \item Number of evaluation runs: yes, \S\ref{sec:setup}; the search log records every evaluated chain and every skipped extension.
  \item Definition of each reported statistic: yes, \S\ref{sec:metrics}.
  \item Central tendency and variation: yes; prompt-clustered bootstrap intervals for rates, document-level Wilson intervals beside them, mean and standard deviation elsewhere.
  \item Runtime or energy cost: \ph{TODO}.
  \item Computing infrastructure: \ph{TODO}.
\end{itemize}
